\section{Introduction}
\lblsec{intro}

What did Claude Monet see as he placed his easel by the bank of the Seine near Argenteuil on a lovely spring day in 1873 (\reffig{teaser}, top-left)? A color photograph, had it been invented, may have documented a crisp blue sky and a glassy river reflecting it. Monet conveyed his {\em impression} of this same scene through wispy brush strokes and a bright palette.\blfootnote{* indicates equal contribution}

What if Monet had happened upon the little harbor in Cassis on a cool summer evening (\reffig{teaser}, bottom-left)? A brief stroll through a gallery of Monet paintings makes it possible to imagine how he would have rendered the scene: perhaps in pastel shades, with abrupt dabs of paint, and a somewhat flattened dynamic range.



\begin{figure}
 \centering
 \includegraphics[width=1.0\hsize]{figs/paired_unpaired2.pdf}
 \vspace{-0.3in}
  \caption{\emph{Paired} training data (left) consists of training examples $\{x_i,y_i\}_{i=1}^N$, where the correspondence between $x_i$ and $y_i$ exists~\cite{isola2016image}. We instead consider \emph{unpaired} training data (right), consisting of a source set $\{x_i\}_{i=1}^N$ ($x_i \in X$) and a target set $\{y_j\}_{j=1}^M$ ($y_j \in Y$), with no information provided as to which $x_i$ matches which $y_j$.}
 \lblfig{paired_unpaired}
 \vspace{-0.2in}
\end{figure}

We can imagine all this despite never having seen a side by side example of a Monet painting next to a photo of the scene he painted. Instead, we have knowledge of the set of Monet paintings and of the set of landscape photographs. We can reason about the stylistic differences between these two sets, and thereby imagine what a scene might look like if we were to ``translate'' it from one set into the other.

In this paper, we present a method that can learn to do the same: capturing special characteristics of one image collection and figuring out how these characteristics could be translated into the other image collection, all in the absence of any paired training examples. 

This problem can be more broadly described as image-to-image translation~\cite{isola2016image}, converting an image from one representation of a given scene, $x$, to another, $y$, e.g., grayscale to color, image to semantic labels, edge-map to photograph.
Years of research in computer vision, image processing, computational photography, and graphics have produced powerful translation systems in the supervised setting, where example image pairs $\{x_i,y_i\}_{i=1}^N$ are available (\reffig{paired_unpaired}, left), e.g., 
\cite{eigen2015predicting, hertzmann2001image, isola2016image, johnson2016perceptual, laffont2014transient, long2015fully, shih2013data, wang2016generative, xie2015holistically, zhang2016colorful}.
However, obtaining paired training data can be difficult and expensive. For example, only a couple of datasets exist for tasks like semantic segmentation (e.g.,~\cite{Cordts2016Cityscapes}), and they are relatively small.  
Obtaining input-output pairs for graphics tasks like artistic stylization can be even more difficult since the desired output is highly complex, typically requiring artistic authoring. For many tasks, like object transfiguration (e.g., zebra$\leftrightarrow$horse, ~\reffig{teaser} top-middle), the desired output is not even well-defined.


We therefore seek an algorithm that can learn to translate between domains without paired input-output examples (\reffig{paired_unpaired}, right). We assume there is some underlying relationship between the domains -- for example, that they are two different renderings of the same underlying scene -- and seek to learn that relationship. Although we lack supervision in the form of paired examples, we can exploit supervision at the level of sets: we are given one set of images in domain $X$ and a different set in domain $Y$. We may train a mapping $G: X \rightarrow Y$ such that the output $\hat{y} = G(x)$, $x \in X$, is indistinguishable from images $y \in Y$ by an adversary trained to classify $\hat{y}$ apart from $y$. In theory, this objective can induce an output distribution over $\hat{y}$ that matches the empirical distribution $p_{data}(y)$ (in general, this requires $G$ to be stochastic)~\cite{goodfellow2014generative}. The optimal $G$ thereby translates the domain $X$ to a domain $\hat{Y}$ distributed identically to $Y$. However, such a translation does not guarantee that an individual input $x$ and output $y$ are paired up in a meaningful way -- there are infinitely many mappings $G$ that will induce the same distribution over $\hat{y}$. Moreover, in practice, we have found it difficult to optimize the adversarial objective in isolation: standard procedures often lead to the well-known problem of mode collapse, where all input images map to the same output image and the optimization fails to make progress~\cite{goodfellow2016nips}.

These issues call for adding more structure to our objective. Therefore, we exploit the property that translation should be ``cycle consistent", in the sense that if we translate, e.g., a sentence from English to French, and then translate it back from French to English, we should arrive back at the original sentence~\cite{brislin1970back}. Mathematically, if we have a translator $G: X \rightarrow Y$ and another translator $F: Y \rightarrow X$, then $G$ and $F$ should be inverses of each other, and both mappings should be bijections. We apply this structural assumption by training both the mapping $G$ and $F$ simultaneously, and adding a \emph{cycle consistency loss}~\cite{zhou2016learning} that encourages $F(G(x)) \approx x$ and $G(F(y)) \approx y$. Combining this loss with adversarial losses on domains $X$ and $Y$ yields our full objective for unpaired image-to-image translation.


We apply our method to a wide range of applications, including collection style transfer, object transfiguration, season transfer and photo enhancement. We also compare against previous approaches that rely either on hand-defined factorizations of style and content, or on shared embedding functions, and show that our method outperforms these baselines. We provide both \href{https://github.com/junyanz/pytorch-CycleGAN-and-pix2pix}{PyTorch} and \href{https://github.com/junyanz/CycleGAN}{Torch} implementations. Check out more results at our \href{https://junyanz.github.io/CycleGAN/}{website}.



\begin{figure*}[th]
 \centering
 \includegraphics[width=1.0\hsize]{figs/system_diagram_v2.pdf}
 \vspace{-0.2in}
  \caption{(a) Our model contains two mapping functions $G: X \rightarrow Y$ and $F: Y \rightarrow X$, and associated adversarial discriminators $D_Y$ and $D_X$. $D_Y$ encourages $G$ to translate $X$ into outputs indistinguishable from domain $Y$, and vice versa for $D_X$ and $F$. To further regularize the mappings, we introduce two {\it cycle consistency losses} that capture the intuition that if we translate from one domain to the other and back again we should arrive at where we started: (b) forward cycle-consistency loss: $x\rightarrow G(x) \rightarrow F(G(x)) \approx x$, and (c) backward cycle-consistency loss: $y \rightarrow F(y) \rightarrow G(F(y)) \approx y$ }
 \lblfig{overview}
 \vspace{-0.2in}
\end{figure*}

